\chapter{Thảo luận, kết luận và hướng phát triển}

\section{Đóng góp của đồ án}
Đồ án cụ thể hóa bài toán ánh xạ mã thành các đối tượng, quan hệ, phép
xử lý và luật có thể thực thi. Điều quan trọng đối với môn Biểu diễn tri
thức là mối liên hệ giữa ba phần: tri thức được lưu dưới dạng gì, bộ suy
luận sử dụng phần tri thức nào và người dùng kiểm tra kết quả bằng cách nào.
Các module và công thức trong báo cáo thể hiện trực tiếp mối liên hệ đó.

Kết quả thứ nhất là hợp đồng dữ liệu phân biệt assertion, assessment và
proof. Kết quả thứ hai là bộ suy luận có thứ tự, có đầu ra từ chối rõ ràng.
Kết quả thứ ba là bộ thực nghiệm vừa báo cáo độ chính xác, vừa thể hiện
phần tri thức chưa được bao phủ. Những kết quả này giúp việc phát triển
thêm luật có điểm xuất phát đo được.

\section{Ý nghĩa trong một quy trình sử dụng}
Trong một quy trình rà soát mapping, hệ thống có thể hỗ trợ người dùng
nhóm các cặp theo assessment và mở bằng chứng để kiểm tra.
Một chính sách nghiệp vụ tương lai có thể được mô hình hóa là
\begin{equation}
 d(m)=\operatorname{Policy}\bigl(f_{\K}(m),\operatorname{context}(m)\bigr).
\end{equation}
Ở đây $d$ là quyết định phê duyệt hoặc yêu cầu xem xét.
\src{Policy} và \src{context} chưa được cài đặt trong repo.
Việc tách này ngăn diễn giải sai rằng một mức A/B/C tự động tương đương
với quyết định nghiệp vụ cuối cùng.

Proof trace cho phép kiểm tra và phát hiện sai sót ở từng luật.
\Unassessed{} tạo ra danh sách cụ thể các ca cần bổ sung tri thức.
Đây là cơ sở để xây dựng công cụ hỗ trợ chuyên gia và theo dõi chất lượng
dữ liệu. Khả năng giảm thời gian review hoặc giảm chi phí cần được đánh
giá bằng một quy trình có người dùng và tiêu chí chấp nhận xác định.

\section{Giới hạn cần nhìn nhận}
\begin{table}[htbp]
\centering\small
\caption{Khoảng cách giữa hiện thực hiện tại và hệ thống hoàn chỉnh}\label{tab:gaps}
\begin{tabularx}{\textwidth}{p{3.0cm}YY}
\toprule
Khía cạnh & Hiện tại & Phần cần bổ sung\\
\midrule
Ngữ nghĩa & Chuỗi, bí danh nhỏ, tập từ, cặp từ đối lập. &
Đồng nghĩa có nguồn, phủ định, quan hệ phân cấp có kiểm chứng.\\
COKB & Sáu nhóm thành phần được vận dụng trong miền hẹp. &
Đối tượng có hành vi nội tại và bộ giải tổng quát nếu mục tiêu mở rộng đòi hỏi.\\
Provenance & ID nguồn và mã luật. &
Phiên bản nguồn, hash luật, thời điểm, tác nhân và lịch sử phê duyệt.\\
Mapping phức hợp & Target đơn trong Python. &
Ngữ nghĩa OneOf/AllOf và kiểm thử tương ứng.\\
Đánh giá & 500 cặp gold, coverage thấp. &
Bộ holdout độc lập, ablation, kiểm định với chuyên gia.\\
Triển khai & CLI và file JSON/Turtle. &
API, quản lý người dùng, policy và quy trình review.\\
\bottomrule
\end{tabularx}
\end{table}

Một hạn chế bổ sung là ontology và rule catalog chưa được dùng làm
nguồn thực thi thống nhất. Thay đổi Turtle hoặc JSON catalog mà không sửa
Python có thể làm tài liệu và hành vi khác nhau. Tính nhất quán giữa các
biểu diễn cần được kiểm tra khi mở rộng.

Graph đầy đủ chưa có trong checkout. Indexer có cơ chế nạp dữ liệu mẫu
nếu không đọc được TTL thực, nên việc một lệnh indexing hoàn thành
không đủ chứng minh rằng toàn bộ dữ liệu ICD đã được nạp.
Trong lần kiểm tra dùng cho báo cáo, store có 0 quad. Cần kiểm tra số lượng,
hệ mã và provenance sau import trước khi thực nghiệm quy mô lớn.

\section{Hướng phát triển ưu tiên}
Ưu tiên thứ nhất là nâng chất lượng tri thức nhãn: bổ sung xử lý phủ định
và đồng nghĩa có kiểm thử hồi quy, sau đó đánh giá trên tập độc lập.
Ưu tiên thứ hai là hoàn thiện provenance và proof: lưu phiên bản nguồn,
hash ruleset, các tập từ đã tính và lý do các luật không khớp.
Ưu tiên thứ ba là phục hồi dữ liệu graph đầy đủ, xử lý đúng hệ mã của từng
dòng và báo cáo số lượng mapping thực sự đã nạp.

Sau các bước đó có thể thử tích hợp lớp giao tiếp ngôn ngữ tự nhiên.
Kết quả do model tạo cần có loại assessment riêng và được đánh giá đối
chiếu với luật; hiện chưa có thực nghiệm hybrid trong đồ án.
Nếu tiếp tục theo COKB tổng quát, cần thiết kế rõ tập loại sự kiện,
hành vi đối tượng và chiến lược suy diễn, thay vì chỉ bổ sung thêm tên
class vào ontology.

\section{Kết luận}
Đồ án cho thấy cách đưa tri thức kiểm tra mapping vào một cấu trúc có
thể vận hành: đối tượng và quan hệ để biểu diễn, toán tử và luật để
đánh giá, proof để giải thích, validation để kiểm tra dữ liệu.
Thực nghiệm xác nhận các thành phần này hoạt động trên phạm vi hiện có,
đồng thời chỉ ra giới hạn bao phủ 4,40\%.
Ý nghĩa thực tế của COKB ở đây là cung cấp những nhận định có thể kiểm tra
và nhận diện rõ các ca chưa đủ tri thức, làm nền cho quy trình rà soát
mapping có sự tham gia của chuyên gia.
